\documentclass[12pt]{article}
\usepackage[T1]{fontenc}
\usepackage[utf8]{inputenc}
\usepackage{lmodern}
\usepackage{textcomp}
\usepackage{enumitem}
\usepackage{geometry}
\geometry{margin=1in}
\begin{document}
\begin{center}
Answer Key - Philosophy - Graduate Level Assessment 7 \
\small (Answers are ordered exactly as in the question paper.)
\end{center}

\section*{Multiple Choice}
\begin{enumerate}[label=\arabic*.]
\item \textbf{Answer:} A
\\ \textit{Options:} A) agent-relative.; B) agent-neutral.; C) absolute.; D) none of the above.
\\ \textit{Note:} Parfit argues that the obligation to prioritize one's children's welfare is agent-relative because it is based on the special personal relationships and responsibilities that exist between a parent and their children, which create a unique moral duty that differs from impartial considerations of welfare.
\item \textbf{Answer:} A
\\ \textit{Options:} A) justice.; B) supererogation.; C) honesty.; D) morality.; E) virtue.
\\ \textit{Note:} Mill argues that utilitarianism can lead to injustices by prioritizing overall happiness over individual rights, thereby neglecting the moral importance of fairness and justice in ethical considerations.
\item \textbf{Answer:} A
\\ \textit{Options:} A) Sri Lakshmi; B) Brahma; C) Nanak; D) Vishnu; E) Shiva
\\ \textit{Note:} Sri Lakshmi is regarded as a central figure in the origin of sacred literature due to her embodiment of wisdom and spiritual knowledge, which are foundational to the philosophical teachings that shaped early religious texts.
\item \textbf{Answer:} C
\\ \textit{Options:} A) TRUE; B) undeniable; C) FALSE; D) unchangeable; E) an illusion
\\ \textit{Note:} Descartes initially doubts the veracity of everything he perceives in order to establish a foundation for certain knowledge, leading him to conclude that what he sees may be false.
\item \textbf{Answer:} A
\\ \textit{Options:} A) in a utilitarian way.; B) in a Cartesian way.; C) in a nihilistic way.; D) in a Hegelian way.; E) in a Kantian way.
\\ \textit{Note:} Nussbaum argues that individuals from diverse conceptual frameworks often prioritize practical outcomes and benefits in their interactions, reflecting a utilitarian perspective that emphasizes efficiency and utility in cross-cultural communication.
\end{enumerate}

\section*{True / False}
\begin{enumerate}[label=\arabic*.]
\setcounter{enumi}{5}
\item \textbf{Answer:} True
\\ \textit{Options:} A) True; B) False
\\ \textit{Note:} Hume posits that moral judgments arise from human emotions and social experiences, asserting that cultural norms significantly influence our understanding of morality.
\item \textbf{Answer:} False
\\ \textit{Options:} A) True; B) False
\\ \textit{Note:} Arthur's assertion is false because Singer explicitly addresses principles of equality and harm in his discussions on world hunger and poverty, advocating for a moral obligation to alleviate suffering.
\item \textbf{Answer:} True
\\ \textit{Options:} A) True; B) False
\\ \textit{Note:} Pence contends that critics of somatic cell nuclear transfer (SCNT) often presuppose that parental objections stem from a sincere desire to protect their child's welfare, thus reflecting a belief in the authenticity of parental motivation.
\item \textbf{Answer:} True
\\ \textit{Options:} A) True; B) False
\\ \textit{Note:} Papadaki aligns with Kant's ethical framework by emphasizing the importance of treating individuals as ends in themselves rather than as means to an end, thereby supporting Kant's stance against sexual objectification.
\item \textbf{Answer:} False
\\ \textit{Options:} A) True; B) False
\\ \textit{Note:} The statement is false because Caityavasis are actually lay followers who reside in rural areas and are characterized by their devotion to stupas and pilgrimage sites, rather than primarily living in urban centers to support monastic communities.
\end{enumerate}

\section*{Long Form Response}
\begin{enumerate}[label=\arabic*.]
\setcounter{enumi}{10}
\item \textbf{Answer:} In Aquinas's natural law theory, intrinsic goods are fundamental values that are inherently valuable and contribute to human flourishing. These goods, such as life, knowledge, and friendship, are significant because they guide moral action and establish a framework for discerning what is good. Aquinas argues that these intrinsic goods are interrelated; for instance, the pursuit of knowledge enhances one's understanding of friendship and life itself. This interconnectedness underscores the holistic nature of Aquinas's ethical framework, as the realization of one intrinsic good often supports the attainment of others, ultimately aligning human actions with divine law. Thus, intrinsic goods serve as both the foundation and a guiding compass within Aquinas's moral philosophy, illustrating the profound relationship between human nature and ethical conduct.
\\ \textit{Note:} correct because it accurately identifies intrinsic goods in Aquinas's natural law theory as essential values that promote human flourishing, emphasizes their role in guiding moral action, and highlights their interrelated nature, which illustrates the holistic approach of Aquinas's ethical framework.
\item \textbf{Answer:} Feinberg's perspective on human desires, particularly in the context of hunger, challenges traditional classifications by arguing that such desires cannot be neatly categorized as either intrinsic or extrinsic. He posits that hunger reflects a complex interplay of biological, psychological, and social factors, which defy simple classification. This complexity suggests that desires are not solely motivated by internal needs or external circumstances, but rather by a dynamic relationship between the two. Consequently, Feinberg emphasizes the inadequacy of conventional frameworks in fully capturing the nuanced nature of human desires, particularly as they relate to fundamental needs like hunger. Thus, he concludes that neither traditional intrinsic nor extrinsic classifications can encompass the depth of human motivation.
\\ \textit{Note:} correct because it captures Feinberg's argument that hunger is a multifaceted desire influenced by various intertwined factors, which complicates its classification as either intrinsic or extrinsic.
\end{enumerate}

\vfill
\noindent\textit{End of Answer Key}
\end{document}